\section{Doubled-spin conventions and the algebraic carrier class}
\subsection{Doubled spins and channels}
All representation labels are nonnegative integers representing twice the ordinary spin.  Thus $a\in\N_0$ corresponds to spin $a/2$.  For doubled labels $a,b$, define
\begin{equation}
\mathcal C(a,b)=\{|a-b|,|a-b|+2,\ldots,a+b\}.
\label{eq:pairchannels}
\end{equation}

\begin{definition}[Ordered carrier]
An ordered carrier is a four-tuple
\[
J=(a,b,c,d)\in\N^4.
\]
Its two channel sets are
\begin{equation}
E(J)=\mathcal C(a,b)\cap\mathcal C(c,d),
\qquad
G(J)=\mathcal C(b,c)\cap\mathcal C(a,d).
\label{eq:channelsets}
\end{equation}
\end{definition}

\begin{definition}[Algebraic two-channel class]
\label{def:algebraicclass}
The class $\mathfrak A_6$ consists of ordered carriers $J=(a,b,c,d)$ satisfying:
\begin{enumerate}[label=(\alph*)]
\item $1\le a,b,c,d\le6$ and $a+b+c+d$ is even;
\item $|E(J)|=|G(J)|=2$;
\item for every $e\in E(J)$ and $f\in G(J)$, the finite Racah interval defined in \eqref{eq:qracah} is nonempty;
\item every triangle, numerator, denominator and dimension-amplitude q-index occurring in \eqref{eq:qdelta}--\eqref{eq:rawF} lies in $\{0,1,\ldots,10\}$.
\end{enumerate}
No floating-point orthogonality tolerance or numerical nonzero-speed test enters this definition.
\end{definition}

Exact integer enumeration gives
\begin{equation}
|\mathfrak A_6|=283
\end{equation}
ordered carriers in 24 unordered external-label families.  The full machine-readable list accompanies the manuscript.

\begin{proposition}[Critical q-index ledger]
\label{prop:indexledger}
Across $\mathfrak A_6$, the exact maxima of the q-indices are:
\begin{center}
\begin{tabular}{lc}
\toprule
Occurrence & maximum index\\
\midrule
triangle-factor q-factorials & 10\\
Racah numerator $[z+1]_q!$ & 10\\
Racah denominator q-factorials & 6\\
dimension factors $[e+1]_q,[f+1]_q$ & 8\\
\bottomrule
\end{tabular}
\end{center}
Hence the critical set is $\mathcal N_{\mathrm{crit}}=\{1,\ldots,10\}$.
\end{proposition}

\begin{proof}
This is an exact finite integer enumeration using Definition \ref{def:algebraicclass}; no transcendental or floating-point calculation is involved.  Pseudocode and the 283-carrier list are given in Appendix \ref{app:pseudocode} and the supplementary JSON file.
\end{proof}

\section{Trigonometric q-Racah convention}
Let
\begin{equation}
q=e^{i\theta},
\qquad
\qnum{n}=\frac{\sin(n\theta)}{\sin\theta},
\qquad
\qfac{n}=\prod_{k=1}^{n}\qnum{k},
\qquad
\qfac{0}=1.
\label{eq:qnumber}
\end{equation}
The anchor is
\begin{equation}
\theta_0=\frac{\pi}{12}.
\label{eq:anchor}
\end{equation}
At this specialization $[12]_q=0$.  Proposition \ref{prop:indexledger} shows that the declared model uses only indices at most ten.  We therefore work with the explicitly defined finite trigonometric q-Racah expression below and do not import an unspoken claim about a complete modular tensor category at the root of unity.

For an admissible triple $(a,b,e)$, write
\begin{align}
\alpha&=\frac{a+b-e}{2},&
\beta&=\frac{a-b+e}{2},&
\gamma&=\frac{-a+b+e}{2},&
\sigma&=\frac{a+b+e}{2}+1.
\end{align}
Define
\begin{equation}
\Delta_q(a,b,e)=
\left(
\frac{\qfac{\alpha}\qfac{\beta}\qfac{\gamma}}
{\qfac{\sigma}}
\right)^{1/2}.
\label{eq:qdelta}
\end{equation}
For six admissible doubled labels, set
\begin{align}
A_1&=\frac{a+b+e}{2},& A_2&=\frac{a+d+f}{2},&
A_3&=\frac{c+b+f}{2},& A_4&=\frac{c+d+e}{2},\\
B_1&=\frac{a+b+c+d}{2},&
B_2&=\frac{a+c+e+f}{2},&
B_3&=\frac{b+d+e+f}{2}.
\end{align}

\begin{definition}[Scalar q-$6j$ convention]
\label{def:q6jconvention}
Throughout this paper,
\begin{align}
\sixj{a}{b}{e}{c}{d}{f}
&=
\Delta_q(a,b,e)\Delta_q(a,d,f)
\Delta_q(c,b,f)\Delta_q(c,d,e)
\nonumber\\
&\quad\times
\sum_{z=\max A_i}^{\min B_j}
(-1)^z
\frac{\qfac{z+1}}
{\prod_{i=1}^{4}\qfac{z-A_i}\prod_{j=1}^{3}\qfac{B_j-z}}.
\label{eq:qracah}
\end{align}
Equation \eqref{eq:qracah}, together with the anchor-selected square-root branches, is the definition used by the manuscript and software.  No asymptotic replacement is used.
\end{definition}

This is the real trigonometric normalization underlying the recoupling conventions used in \cite{kirillovReshetikhin1989,kauffmanLins1994,taylorWoodward2005}; the explicit definition above takes precedence over convention-dependent terminology.

\section{Two-channel recoupling matrix and analytic branches}
For $E(J)=\{e_1,e_2\}$ and $G(J)=\{f_1,f_2\}$ in increasing order, define
\begin{equation}
\widetilde F_{ij}(\theta)=
(-1)^{(a+b+c+d)/2}
\sqrt{\qnum{e_i+1}\qnum{f_j+1}}
\sixj{a}{b}{e_i}{c}{d}{f_j}.
\label{eq:rawF}
\end{equation}
The phase and dimension factors in \eqref{eq:rawF} are part of the convention.  A constant signed diagonal gauge $G_J$ is chosen at $\theta_0$ to give positive orientation and a fixed sign for the $(1,1)$ entry:
\begin{equation}
F_J(\theta)=G_J\widetilde F_J(\theta).
\label{eq:gaugedF}
\end{equation}
The gauge is frozen as $\theta$ varies.

\begin{lemma}[Holomorphic square-root continuation]
\label{lem:sqrtbranch}
Let $D\subset\C$ be simply connected, let $r:D\to\C\setminus\{0\}$ be holomorphic, and fix one square root of $r(z_0)$ at $z_0\in D$.  There is a unique holomorphic square root on $D$ with that anchor value.
\end{lemma}

\begin{proof}
A nonvanishing holomorphic function on a simply connected domain has a holomorphic logarithm.  Half of that logarithm, followed by exponentiation, gives the required branch; the anchor value fixes the sign.
\end{proof}

Lemma \ref{lem:sqrtbranch} applies both to the four triangle radicands and to the dimension-amplitude radicands in \eqref{eq:rawF} once their q-factorials are separated from zero.

\begin{proposition}[Transpose orthogonality]
\label{prop:orthogonality}
For the convention \eqref{eq:qracah}--\eqref{eq:rawF}, the standard q-$6j$ orthogonality identity gives
\begin{equation}
F_J(\theta)F_J(\theta)^T=I_2
\label{eq:orthogonality}
\end{equation}
for real $\theta$ in every regular chamber.  On a connected complex regular disk, \eqref{eq:orthogonality} continues as an analytic transpose-orthogonality identity.  It is not a unitarity assertion.
\end{proposition}

\begin{proof}
The dimension-weighted q-$6j$ orthogonality relation is the change-of-basis identity for the two complete channel sets; see \cite{kirillovReshetikhin1989,kauffmanLins1994}.  The constant signed gauge preserves it.  Both sides are analytic on a regular disk and agree on a real interval, so the identity theorem gives the complex continuation.
\end{proof}

Define
\begin{equation}
K_J(\theta)=F'_J(\theta)F_J(\theta)^T.
\label{eq:generator}
\end{equation}
Differentiating \eqref{eq:orthogonality} yields $K_J^T=-K_J$.  Hence
\begin{equation}
K_J(\theta)=
\begin{pmatrix}
0&-\omega_J(\theta)\\
\omega_J(\theta)&0
\end{pmatrix}.
\label{eq:omega}
\end{equation}
The scalar $\omega_J$ is the local angular speed of the recoupling frame.
